%! Summary Statistics for a Quantitative Variable — Unit 1, Lesson 1.5 — Homework
% EFFL, AP Statistics variant: this is the lesson's SCORED practice and the LAST component in
% the packet. These items were the in-class Check Your Understanding set; they are now the
% graded homework, started in the closing minutes of the period and finished at home. The
% AP Practice component that precedes it stays assigned-but-unscored.
% \answerspace{H}{} reserves exactly H in the blank; the key passes the answer as the second
% argument, so the two files paginate identically by construction. Keep the macro and every
% \boxguard byte-identical in the blank and the key.
%
% ANSWERS FOR REFERENCE (the key carries the prose):
%   1. 16 12 60 15 19 14 18 -> sum 154, mean 22 min; ordered 12 14 15 16 18 19 60, median 16.
%   2. 8 9 10 11 12 (mean 10, median 10, range 4); 12 -> 32 gives mean 14, median 10, range 24.
%   3. 15 25 30 40 45 55 | 60 70 75 85 100 120 -> median 57.5, Q1 35, Q3 80, IQR 45.
%   5. (B) — mean far above median means right skew from a few very large salaries.
\documentclass[10pt]{article}
\usepackage{apstats-article}
\usepackage{apstats-boxes}
\newcommand{\answerspace}[2]{\par\nopagebreak\noindent\begin{minipage}[t][#1][t]{\linewidth}%
  \color{keyred}\bfseries #2\end{minipage}\par}

\begin{document}

\pageheader{Unit 1, Lesson 1.5}{Homework}
% NAMESTRIP: no \namedateperiod here — the name row lives on the cover only.

\vspace{0.15in}

% Authored BYTE-IDENTICALLY in homework/ and homework_key/ — this box is what keeps the
% blank and the key the same number of pages.
\boxguard[8]
\begin{remindbox}
\small
\textbf{This is your graded homework.} It is scored and due next class. You will start it in the
last minutes of the period, so ask about anything that stalls you before you leave, then finish
it on your own.
\end{remindbox}

\vspace{0.12in}

% ── The scored practice (formerly this lesson's Check Your Understanding) ────
\boxguard[14]
\begin{notesbox}{Choosing and Computing a Summary}
Show the arithmetic wherever a question asks for a number, and write every explanation
\emph{in the context of the problem} --- that is what earns credit on the AP exam, and it is the
standard here.
\begin{enumerate}[leftmargin=1.9em, itemsep=6pt]
  \item Seven employees at a small office reported their commute times, in minutes:

        \vspace{3pt}
        \begin{center}
        16 \quad 12 \quad 60 \quad 15 \quad 19 \quad 14 \quad 18
        \end{center}
        \vspace{3pt}

        \textbf{(a)} Mean: \blank{2.6cm} \qquad Median: \blank{2.6cm}

        \vspace{5pt}
        \textbf{(b)} Which of the two better represents a typical commute for this office?
        Use the \emph{shape} of the data to justify your choice.
        \answerspace{2.8cm}{}
  \item Five children reported how many books they read over the summer:
        \quad 8 \quad 9 \quad 10 \quad 11 \quad 12. \quad For this set the mean is 10, the
        median is 10, and the range is 4.

        \vspace{4pt}
        The librarian then discovers a recording error: the child credited with 12 books
        actually read \textbf{32}. Recompute:

        \vspace{4pt}
        New mean: \blank{2.2cm} \qquad New median: \blank{2.2cm} \qquad New range:
        \blank{2.2cm}

        \vspace{5pt}
        In one or two sentences, say which summaries moved, which did not, and why.
        \answerspace{2.6cm}{}
  \item Twelve students recorded how many minutes they spent on homework one evening. The
        values are already in order:

        \vspace{3pt}
        \begin{center}
        15 \quad 25 \quad 30 \quad 40 \quad 45 \quad 55 \quad 60 \quad 70 \quad 75 \quad
        85 \quad 100 \quad 120
        \end{center}
        \vspace{3pt}

        Median: \blank{2.0cm} \quad $Q_1$: \blank{2.0cm} \quad $Q_3$: \blank{2.0cm} \quad
        IQR: \blank{2.0cm}

        \vspace{5pt}
        Show how you split the list to get $Q_1$ and $Q_3$.
        \answerspace{2.4cm}{}
  \item A pediatrician's growth chart reports that the \textbf{90th percentile} of birth weight
        for newborn girls is 3{,}900 grams. Write one sentence interpreting this value in
        context.
        \answerspace{2.4cm}{}
  \item \textbf{AP-style multiple choice.} A company with 40 employees reports that the
        \emph{mean} salary is \$95{,}000 while the \emph{median} salary is \$58{,}000. Which
        statement is the best conclusion? Circle one, then justify it in a sentence.
        \begin{enumerate}[label=\textbf{(\Alph*)}, leftmargin=2.2em, itemsep=2pt]
          \item Half of the employees earn more than \$95{,}000.
          \item The distribution of salaries is likely skewed to the right, with a few very high
                salaries pulling the mean above the median.
          \item An error was made, because the mean and the median of a data set must be equal.
          \item The distribution of salaries is likely skewed to the left.
          \item The standard deviation of the salaries must be zero.
        \end{enumerate}
        \answerspace{2.6cm}{}
\end{enumerate}
\end{notesbox}

\vspace{0.12in}

% The next-lesson preview closes the packet, so it lives here rather than in AP Practice.
\boxguard[10]
\begin{spiralbox}
\small
\textbf{Next lesson: Boxplots and Comparing Distributions.} You can now pull five numbers out of
a data set --- the minimum, $Q_1$, the median, $Q_3$, and the maximum. Next time those five
numbers become a picture, and that picture is what lets us put two groups side by side and say,
with evidence, how they differ. The $1.5\times\text{IQR}$ rule finally goes to work there too:
it is what decides where the picture's whiskers stop.
\end{spiralbox}

\end{document}
